The main text's \S\ref{sec:setup} states what this bound did and did not do
for the paper. This appendix gives the formal development: the coding model,
the hard distribution, and the two theorems. Proofs are in
Appendix~\ref{app:proofs}.

\subsection{Rate-\texorpdfstring{$R$}{R} merging codes}
\label{app:rd-codes}

A rate-$R$ merging code is a pair $(E, D)$ with encoder
$E : (\R^d)^T \to \{1,\ldots,2^R\}$ and decoder
$D : \{1,\ldots,2^R\} \to \R^d$, both of which may depend on
$\boldsymbol H$ as shared side information: in deployment $\boldsymbol H$ is
determined by the LoRA factors both ends already hold, so only the $R$ bits
describing $\boldsymbol\tau$ are charged to the rate. Our lower bound holds
even when $\boldsymbol H$ is revealed for free to both ends. The
\emph{rank-$r$ rate-distortion function} is
\begin{equation}\label{eq:rd-def}
  \mathcal{D}^\star(T, d, r, B, R)
  \;:=\; \inf_{(E, D)}\;\sup_{(\boldsymbol\tau, \boldsymbol H)
    \text{ admissible}}\;
  \E\bigl[\Delta(w^\star;\boldsymbol\tau, \boldsymbol H)\bigr],
  \qquad w^\star := D(E(\boldsymbol\tau)),
\end{equation}
with the expectation over any shared randomness, and $\Delta$ the worst-task
distortion \eqref{eq:worst-task}. In the floor-zero Stiefel regime this
appendix determines $\mathcal{D}^\star$ up to a universal constant: a lower
bound for every code (Theorem~\ref{thm:general}) and an explicit code matching
it (Theorem~\ref{thm:achv}).

\paragraph{Hard distribution $P^\star$.} The lower bound is established
against a distribution $P^\star$ on admissible tuples. For each task
$t$, independently: (i) sample $U_t \in \mathrm{Stiefel}(d, r)$
uniformly and set $V_t := \range(U_t)$; (ii) fix a positive-definite
diagonal $D_t \succ 0$ (the task-$t$ spectrum) and set $H_t := U_t D_t
U_t^\top$; (iii) sample $u_t \sim \mathrm{Unif}(S^{r-1})$ and set
$\tau_t := U_t\, B D_t^{-1/2} u_t$. The normalisation
$\tau_t^\top H_t \tau_t = B^2$ holds by construction, and the
$D_t^{-1/2}$ stretch makes the second-moment identity
$\E[H_t \tau_t \tau_t^\top H_t \mid U_t] = (B^2/r) H_t$ hold
\emph{regardless of} $D_t$, the feature that makes
Theorem~\ref{thm:general} apply to arbitrary task-dependent spectra.

\paragraph{Yao's minimax bridge.} By Yao's minimax
principle~\citep{yao1977}, $\mathcal{D}^\star$ is at least the
Bayesian rate-distortion function under the prior $P^\star$:
\begin{equation}\label{eq:yao}
  \mathcal{D}^\star(T,d,r,B,R)
  \;\geq\;
  \inf_{(E,D)}\;\E_{(\boldsymbol\tau,\boldsymbol H)\sim P^\star}
    \bigl[\Delta(w^\star;\boldsymbol\tau,\boldsymbol H)\bigr]
\end{equation}
All lower bounds below are proved against the right-hand side.

\subsection{The lower bound}
\label{app:rd-lower}

The lower bound rests on three ingredients: a max-$\geq$-avg identity
converting max-distortion into average distortion plus an irreducible
floor; a closed-form evaluation of that floor under $P^\star$; and a
Shannon converse on the remaining compression subproblem, conditional
on the curvature data $\boldsymbol H$. Lemmas~\ref{lem:id}
and~\ref{lem:floor-general} are proved in full in
Appendix~\ref{app:proofs}. Theorem~\ref{thm:general}'s floor term is proved
there unconditionally; its rate term rests on Assumption~\ref{asm:epr} below,
which we state explicitly rather than leave in an appendix.

\begin{lemma}[$H_t$-weighted max $\geq$ avg identity]\label{lem:id}
For every admissible $(\boldsymbol\tau, \boldsymbol H)$ and every
$w\in\R^d$,
\begin{equation}\label{eq:id-eq}
  \tfrac1T\sum_{t=1}^T (w-\tau_t)^\top H_t (w-\tau_t)
  \;=\; (w-\tauH)^\top\bar H(w-\tauH) + \tfrac1T\sum_{t=1}^T
  (\tau_t-\tauH)^\top H_t (\tau_t-\tauH),
\end{equation}
and hence $\max_t (w-\tau_t)^\top H_t(w-\tau_t)$ is at least the
right-hand side.
\end{lemma}

The proof is a two-line expansion: the cross-term vanishes by the
defining property $\bar H\tauH = T^{-1}\sum_t H_t\tau_t$ of the
centroid. Two consequences drive everything that follows: components
of $w$ orthogonal to $\range(\bar H) = V_1+\cdots+V_T$ annihilate
against every $H_t$, so without loss the decoder's output lies in an
at-most-$\deff$-dimensional subspace; and the bound separates into a
compression term in the $\bar H$-metric on $\tauH$ and a floor term
that does not depend on $w$ at all.

\begin{lemma}[Closed-form floor]\label{lem:floor-general}
Under $P^\star$ with any fixed positive-definite spectra
$(D_t)_{t=1}^T$ (possibly distinct across tasks),
\begin{equation}\label{eq:floor-closed}
  \E_{P^\star}\!\left[\tfrac1T\sum_{t=1}^T(\tau_t-\tauH)^\top
   H_t(\tau_t-\tauH)\right]
  \;=\; B^2\!\left(1 - \frac{\E[\deff]}{Tr}\right).
\end{equation}
\end{lemma}

The key probabilistic input is the second-moment identity
$\E[H_t\tau_t\tau_t^\top H_t\mid U_t] = (B^2/r) H_t$, which holds for
any $D_t$ by the $D_t^{-1/2}$ stretch in $P^\star$. Three limits
illuminate the formula: at $r = d$ (full rank) the floor is
$B^2(1 - 1/T)$, the classical averaged-vector-quantisation spread; at
$T = 1$ it is zero, recovering single-source quantisation; and when
all subspaces coincide ($V_t = V$, $\deff = r$) it is again
$B^2(1 - 1/T)$. In the generic LoRA regime ($V_t$ Stiefel-random,
$Tr \leq d$), $\deff = Tr$ almost surely and the floor vanishes.

The rate term needs one ingredient we do not prove, and we state it as an
assumption rather than carry it as a caveat, so that what is conditional is
visible where the result is used.

\begin{assumption}[Uniform entropy power]\label{asm:epr}
There is an absolute constant $c_{\mathrm{TQ}} > 0$, independent of
$(T, r, d)$ and of the spectra $(D_t)_{t=1}^T$, such that the conditional
entropy power of $\eta := \bar H^{1/2}\tauH$ given $\boldsymbol H$ is at least
$c_{\mathrm{TQ}}$ times its conditional variance.
\end{assumption}

For each fixed instance the ratio is finite, because $\eta$ has a density and
its scale is pinned by $\tau_t^\top H_t \tau_t = B^2$. What we do not establish
is that the bound is uniform, and Appendix~\ref{app:proofs} isolates the exact
step where it is used. We import the constant from
TurboQuant~\citep{zandieh2025turboquant}, which supplies one in the
rotation-invariant case; transferring it to our conditional setting is part of
what Assumption~\ref{asm:epr} asserts.

\begin{theorem}[Rank-$r$ RD lower bound, general $H_t$]\label{thm:general}
Under $P^\star$ with any fixed positive-definite spectra
$(D_t)_{t=1}^T$ and $\rank(H_t) = r$ for all $t$, and under
Assumption~\ref{asm:epr} for the rate term,
\begin{equation}\label{eq:thm-lb}
  \mathcal{D}^\star(T, d, r, B, R)
  \;\geq\;
  B^2\!\left(1 - \frac{\E[\deff]}{Tr}\right)
  \;+\; c_{\mathrm{TQ}}\cdot\frac{B^2\E[\deff]}{Tr}
  \cdot 2^{-2R/\E[\deff]}.
\end{equation}
Under Stiefel-random sampling with $Tr\leq d$, $\deff = Tr$ a.s., the
floor vanishes, and
$\mathcal{D}^\star(T,d,r,B,R) \geq c_{\mathrm{TQ}} B^2\,
2^{-2R/(Tr)}$, with $c_{\mathrm{TQ}}$ the universal constant of
TurboQuant~\citep{zandieh2025turboquant} Theorem~3.
\end{theorem}

\paragraph{Proof sketch.} By \eqref{eq:yao} it suffices to lower-bound
the Bayesian risk under $P^\star$. Lemmas~\ref{lem:id}
and~\ref{lem:floor-general} reduce it to the floor plus the
compression term $\E[(w^\star-\tauH)^\top\bar H(w^\star-\tauH)]$; the
decoder may be assumed to emit $w^\star \in \range(\bar H)$, and with
$\eta := \bar H^{1/2}\tauH \in \R^{\deff}$ the subproblem is rate-$R$
quantisation of $\eta$ under squared loss. Given $\boldsymbol H$, the
only remaining randomness is the directions $u_t$, so a conditional
Shannon converse applies to every rate-$R$ code, \emph{including
codebooks adapted to $\boldsymbol H$}: conditioning on
$\boldsymbol H$, rather than averaging over rotations, is what makes
the converse valid against the $\boldsymbol H$-aware decoder of
Appendix~\ref{app:rd-achv}; the constant absorbs the entropy-power-to-variance
ratio of $\eta$, which is $O(1)$ because the scale is pinned by
$\tau_t^\top H_t \tau_t = B^2$. Combining with
$\E[\|\eta\|^2] = B^2\E[\deff]/(Tr)$ from
Lemma~\ref{lem:floor-general} yields \eqref{eq:thm-lb}.

\paragraph{On putting $\E[\deff]$ inside the exponent.} The converse is
applied conditionally on $\boldsymbol H$, so what the argument delivers is
$\E[g(\deff)]$ with $g(x) := x\,2^{-2R/x}$, whereas \eqref{eq:thm-lb} is
stated as $g(\E[\deff])$. The substitution needs $g$ convex, and it is, on
all of $x > 0$: writing $a := 2R\ln 2$,
\[
  g''(x) \;=\; \frac{a^{2}}{x^{3}}\,e^{-a/x} \;>\; 0 ,
\]
so Jensen gives $\E[g(\deff)] \geq g(\E[\deff])$, which is the direction a
lower bound requires. The step is worth spelling out because the natural
misreading fails: the exponential factor \emph{alone}, $x \mapsto
2^{-2R/x}$, has second derivative $(a/x^{3})(a/x - 2)e^{-a/x}$ and is
convex only for $x < R\ln 2$, so it would not license the substitution at
the dimensions we work at. It is the linear prefactor, which comes from
$\E[\|\eta\|^2] = B^2\E[\deff]/(Tr)$, that makes the product convex
throughout. The point is moot in the Stiefel regime, where $\deff = Tr$
almost surely and the expectation is over a constant.

\paragraph{Why the general-$H_t$ statement is the headline.} Real LoRA
fine-tunes have task-specific curvature: the Fisher information at
$\tau_t$ is rank-$r$ PSD with range $V_t$ but generally not a
projector. Theorem~\ref{thm:general} applies with no modeling
assumption on the $D_t$ beyond positive-definiteness, closing the
$\mathrm{MSE}\to\mathrm{CE}$ bridge for the lower-bound half.

\subsection{Achievability}
\label{app:rd-achv}

The lower bound is matched by an explicit algorithm in the generic
regime ($V_t$ Stiefel-random, $Tr \leq d$, so $\deff = Tr$ a.s.\ and
the floor vanishes). The algorithm extends the single-source
construction of TurboQuant~\citep{zandieh2025turboquant} with a
centroid-reduction step via Lemma~\ref{lem:id}.

\paragraph{Algorithm.}\label{subsec:alg}\label{subsec:rd-encoder-impl}
Encoder and decoder share a pseudorandom seed indexing an orthogonal
rotation $O \in O(\deff)$ drawn from Haar measure (implemented by QR of
an iid Gaussian matrix). \emph{Encoder:} (1)~compute $\bar H$, $\tauH$,
and an orthonormal basis $Q \in \R^{d\times\deff}$ of $\range(\bar H)$
with eigenvalues $w_1,\ldots,w_{\deff} > 0$; form $\eta := \bar
H^{1/2}\tauH$ in the $Q$-basis. (2)~Apply the shared rotation
$\tilde\eta := O\eta_Q$. (3)~Uniform-scalar-quantise each coordinate on
$[-c\sigma_{\mathrm{pc}}, c\sigma_{\mathrm{pc}}]$ with $b$ bits, where
$\sigma_{\mathrm{pc}} := B/\sqrt{Tr}$ and $c$ is a clipping constant
(we use $c = 5$); total rate $R = b\deff$. \emph{Decoder:} invert each
step and emit $w^\star := \bar H^{+1/2}\eta_q$ on $\range(\bar H)$. The
centroid $\tauH$ is the sufficient statistic in the floor-zero regime,
the $\bar H^{1/2}$ scaling whitens it, and the random rotation spreads
quantisation error equally across directions.

\begin{theorem}[Matching achievability, floor-zero Stiefel]\label{thm:achv}
Under $P^\star$ with Stiefel-random $V_t$ and $Tr\leq d$, for any fixed
positive-definite $(D_t)_{t=1}^T$, the algorithm above at rate $R$
achieves
\begin{equation}\label{eq:thm-ub}
  \E_{P^\star}\bigl[\max_t(w^\star-\tau_t)^\top H_t(w^\star-
  \tau_t)\bigr]
  \;\leq\; C\cdot B^2 \cdot 2^{-2R/(Tr)},
  \qquad C = Tc^2/3.
\end{equation}
Combined with Theorem~\ref{thm:general}, the Bayesian rate-distortion
function under $P^\star$ is $\Theta\bigl(B^2\, 2^{-2R/(Tr)}\bigr)$ in
this regime, universal in $d$, $r$, $R$, $B$, with the
$T$-dependence of the constant discussed in Remark~\ref{rem:tgap}.
\end{theorem}

\paragraph{Proof sketch.} In the floor-zero regime
Lemma~\ref{lem:floor-general} forces $H_t\tauH = H_t\tau_t$ almost
surely, so worst-task distortion is driven entirely by the quantisation
error $\delta\eta$ on the whitened centroid; the crude bound
$\max_t \leq \sum_t = T\cdot\mathrm{avg}$ gives
$\max_t(\delta w)^\top H_t(\delta w) \leq T\|\delta\eta\|^2$, and
uniform scalar quantisation after the rotation contributes
per-coordinate MSE at most $c^2\sigma_{\mathrm{pc}}^2 2^{-2b}/3$.

\begin{remark}[What the $\Theta$ hides; the $T$-dependence is open]
\label{rem:tgap}
The lower-bound constant is $T$-free while the upper-bound constant is
$C = Tc^2/3$; the entire factor $T$ enters at the crude
$\max_t \leq T\,\mathrm{avg}$ step, which is loose except when one
task's error projection dominates, a case atypical in the floor-zero
regime. The synthetic measurements decide this in favour of the benign
reading: sweeping $T \in \{2,4,8,16\}$ ($1000$ trials/cell,
Appendix~\ref{app:synthetic}), the achievability constant grows
\emph{logarithmically} ($R^2 \geq 0.998$ for a $\log T$ fit vs.\
$\approx 0.91$ linear; the $T{=}16$ constant is $1.77\times$ the
$T{=}2$ value, against $8\times$ if linear). The linear $T$ is an
artifact of the analysis, not intrinsic; proving the matching $\log T$
(or $T$-free) constant is open, and the floor-zero $\Theta$ should be
read as ``$\Theta$ up to a $\log T$ factor.''
\end{remark}

\paragraph{The shared-$V$ / floor-positive regime.}\label{subsec:sharedv}
The regime below is a property of the model, not of any cohort we
measured: \S\ref{subsec:floor-zero} finds $\deff = Tr$ on every real
cohort, in both initialisation conditions, so none of them occupies it.
It is reachable when the $V_t$ genuinely coincide, which for LoRA would
need $Tr > d$ or exactly shared subspaces rather than merely
near-collinear ones. When they do overlap ($\deff < Tr$, floor positive),
the zero-floor step fails. A refinement quantising around the Chebyshev center of the
ensemble in an active-set-aware null-space split achieves excess
distortion of order $2^{-2R/(\deff + |A| - 1)}$, where $|A|$ is the
active-set size: better than naive scalar quantisation but short of
the lower bound's $2^{-2R/\deff}$.

\begin{remark}[Shared-$V$ exponent gap]\label{rem:sharedv-gap}
Valid rate-$R$ random-codebook encoders empirically achieve slopes
strictly steeper than the linear construction (slope $-1.45$ vs.\ the
linear prediction $-1.20$ at $T = 3$, $r = 3$;
Appendix~\ref{app:synthetic}), so the true rate-distortion function in
the floor-positive regime lies strictly between the two bounds. Neither
bound is tight; the exact exponent is open.
\end{remark}
